\documentclass[11pt,a4paper]{article}
% Common preamble for the hint sheets (\input from each file)
\usepackage[utf8]{inputenc}
\usepackage[T1]{fontenc}
\usepackage{lmodern}
\usepackage{amsmath,amssymb,amsthm,mathtools}
\usepackage[a4paper,margin=2.3cm]{geometry}
\usepackage{enumitem}
\usepackage{xcolor}
\usepackage[hidelinks]{hyperref}
\renewcommand{\leq}{\leqslant}
\renewcommand{\le}{\leqslant}
\renewcommand{\geq}{\geqslant}
\renewcommand{\ge}{\geqslant}
\setlength{\parindent}{0pt}
\setlength{\parskip}{0.35em}
% \hintstitle{Sheet title}{Class work / Homework}
\newcommand{\hintstitle}[2]{\begin{center}{\LARGE\bfseries Hints: #1\par}\vspace{0.5em}{\large #2\par}\end{center}
\noindent\textit{The hints to each problem are cumulative. Read one, try again, and open the next only if you are still stuck.}\par\medskip\hrule\medskip}
\newcommand{\problem}[1]{\par\bigskip{\large\bfseries Problem #1}\par\nopagebreak\smallskip}
\newcommand{\subpart}[1]{\par\smallskip\textbf{(#1)}\par\nopagebreak}
\newcommand{\hint}[1]{\par\smallskip\noindent\textcolor{black!70}{\textbf{Hint #1.}}\enspace}
\newcommand{\note}[1]{\par\smallskip\noindent\textit{Note.} #1\par}
\providecommand{\lt}{<}
\providecommand{\gt}{>}
\providecommand{\tightlist}{\setlength{\itemsep}{0pt}\setlength{\parskip}{0pt}}
\emergencystretch=3em

\begin{document}
\hintstitle{Functional and differential equations -- 2. Functional-differential equations}{Class work}

\problem{1}
\hint{1} Put \(a=t+h\), \(b=t-h\): \(f(t+h)-f(t-h)=2hf'(t)\). With \(h=1\), \(f'(t)=\frac{f(t+1)-f(t-1)}2\) is continuous, and in the same way differentiable: by induction \(f\in C^\infty\).
\hint{2} Differentiate the identity twice in \(h\): \(f''(t+h)-f''(t-h)=0\) for all \(t,h\), so \(f''\) is constant. Answer: \(f(x)=ax^2+bx+c\).

\problem{2}
\hint{1} Regularity first: the geometric mean of \(2x\) and \(\frac x2\) is \(x\), so \(f'(x)=\frac{f(2x)-f(x/2)}{3x/2}\) and \(f\in C^\infty\). Put \(a=ts\), \(b=\frac ts\): \(f(ts)-f(\frac ts)=t\bigl(s-\frac1s\bigr)f'(t)\).
\hint{2} Differentiate in \(s\): with \(\psi(x)=xf'(x)\) you get \(\psi(ts)+\psi(\frac ts)=\bigl(s+\frac1s\bigr)\psi(t)\). Differentiate twice more in \(s\) and put \(s=1\): \(t^2\psi''+t\psi'-\psi=0\).
\hint{3} This Euler equation gives \(\psi(t)=At+\frac Bt\), i.e.\ \(f'(x)=A+\frac B{x^2}\). Answer: \(f(x)=Ax+\frac Cx+D\); check: \(f(a)-f(b)=(a-b)\bigl(A-\frac C{ab}\bigr)\).

\problem{3}
\hint{1} Put \(c=f(0)\), \(d=f'(0)\); \(y=0\) gives \((1-d)f(x)=cf'(x)\). If \(c\ne0\), then \(f=ce^{\lambda x}\), and the equation forces \(c\lambda=\frac12\): \(f(x)=\frac{e^{\lambda x}}{2\lambda}\), \(\lambda\ne0\).
\hint{2} If \(c=0\), either \(f\equiv0\) or \(d=1\). Regularity: if \(f(y_0)\ne0\), then \(f'(x)=\frac{f(x+y_0)-f(x)f'(y_0)}{f(y_0)}\) is differentiable, so \(f\in C^\infty\). Differentiate in \(y\) at \(y=0\): once gives \(f''(0)=0\), twice gives \(f''=kf\) with \(k=f'''(0)\).
\hint{3} With \(f(0)=0\), \(f'(0)=1\): \(f(x)=x\), \(\frac{\sin\omega x}\omega\) or \(\frac{\sinh\omega x}\omega\) (\(\omega\gt0\)), and all of them work. Answer: \(f\equiv0\), \(\frac{e^{\lambda x}}{2\lambda}\) (\(\lambda\ne0\)), \(x\), \(\frac{\sin\omega x}\omega\), \(\frac{\sinh\omega x}\omega\).

\problem{4}
\hint{1} Let \(g=f^2\ge0\). Then \(g'=f\) and \((g')^2=g\). On an interval where \(g\gt0\), the derivative \(f=g'\) does not vanish and, having the Darboux property, keeps its sign; so \((\sqrt g)'=\pm\frac12\), i.e.\ \(\sqrt g\) is linear with slope \(\pm\frac12\).
\hint{2} Each component of \(\{g\gt0\}\) is therefore a half-line: a linear function vanishes at most once, so the component cannot have two finite ends. On \((b,\infty)\) we get \(f(x)=\frac{x-b}2\), on \((-\infty,a)\) we get \(f(x)=\frac{x-a}2\).
\hint{3} Answer: for \(-\infty\le a\le b\le+\infty\), \(f(x)=\frac{x-a}2\) for \(x\lt a\), \(f=0\) on \([a,b]\), \(f(x)=\frac{x-b}2\) for \(x\gt b\) (all of these work; \(a=-\infty\), \(b=+\infty\) gives \(f\equiv0\)).

\problem{5}
\hint{1} For \(x\gt1\), \(f'(x)=f(x)-f(x-1)\). Inductively, if \(f\in C^{k-1}\) on \((k-1,\infty)\), the right-hand side is \(C^{k-1}\) on \((k,\infty)\), so \(f\in C^k\) there. More cannot be said: \(f\) on \([0,1]\) is an arbitrary continuous function with \(f(1)=\int_0^1f\).
\hint{2} Substituting \(e^{\lambda x}\) into the integral equation gives \(\int_0^1e^{-\lambda s}\,ds=1\), i.e.\ \(\lambda=1-e^{-\lambda}\). The root \(\lambda=0\) gives the constants. If \(\operatorname{Re}\lambda\ge0\) and \(\lambda\ne0\), then \(\bigl|\int_0^1e^{-\lambda s}ds\bigr|\lt1\); so all other roots have negative real part, and the corresponding modes decay.
\hint{3} Every solution is bounded: \(f(x)\) is the average of \(f\) over \([x-1,x]\), so \(\sup_{[0,x]}|f|=\max_{[0,1]}|f|\). Moreover \(I(x)=\int_{x-1}^x(t-x+1)f(t)\,dt\) is constant (\(I'=f(x)-\int_{x-1}^xf=0\)). By the bounded variation from Definite integrals -- 4, \(f\) has a limit \(L\), and \(I\equiv\frac L2\) gives \(L=2\int_0^1tf(t)\,dt\).
\end{document}
